\scanpage{264}
- \mspage{76} - \margin{IX}

de prolonger ce groupe de Barsotti-Tate en un sous-groupe de Barsotti-Tate
\(F\) de \(A\), qui jouerait alors le rôle joué par la \uline{partie fixe}
\(T_\ell(A)^{\mathrm f}\) de (2,2.3.2) dans le cas où \(A^0\) est
\(\ell\)-divisible. Dans le cas semi-stable, où on a \(\overline A_0=A_0^0\),la
question se résoud par l'affirmative par la méthode de 2.2. Dans le cas général,
lorsque le problème envisagé a une solution, on voit facilement qu'elle est
unique, et que les sous-pro-groupes correspondants de \(T_p(A_K)\) et de
\(T_p(A'_K)\) (qui méritent alors encore le nom de \uline{partie fixe} des
pro-groupes envisagés) satisfont encore au théorème d'orthogonalité 5.2, enfin que
la partie fixe \(T_p(A_K)^{\mathrm f}\) peut encore se caractériser comme indiqué
dans 5.3.

6.2.1. Montrons cependant qu'en général le problème envisagé n'a pas une
solution affirmative, i.e. qu'\uline{il n'est pas possible en général de}
\uline{définir la partie fixe} \(T_p(A)^{\mathrm f}\) (cas de la réduction non
semi-stable). En effet, en vertu d'un théorème de relèvement de SERRE-TATE [30],
lorsque \(S\) est complet, l'existence de la partie fixe équivaut à celle de
relèvements infinitésimaux \(\overline A_n\) de \(\overline A_0\) dans les
\(A_n\), ou encore à celle d'un schéma en groupes lisse \(\overline A\) sur
\(S\), extension d'un schéma abélien par un tore \(T\), de fibre spéciale
\(\overline A_0\), et d'un homomorphisme formel
\(\widehat{\overline A}\longrightarrow\unsure{\widehat A|A}\) prolongeant
l'inclusion sur les fibres spéciales. Donc, dans le cas où \(\overline A_0\)
est un schéma abélien i.e. où \(T_0=0\), cela équivaut à l'existence d'un schéma
abélien \(\overline A\) prolongeant \(\overline A_0\) et d'un homomorphisme
\(\overline A\longrightarrow A\) prolongeant l'inclusion des fibres spéciales.
Il suffira donc de donner un cas où \(T_0=0\), où \(A_0^0\) est non unipotent
et non propre, et où \(A_K\) est simple sur \(K\).
\margin{\figure{\begin{tikzpicture}[scale=.3]
\draw (0,.7)--(.3,.8)--(.15,.65);\draw (.7,.2)--(.9,.1);
\draw (1.4,0)--(1.55,.05);
\end{tikzpicture}}}
\margin{394}
\margin{394}
\margin{263}
